% ==============================================================================
% =================================== Einleitung ================================
% ==============================================================================

\chapter{Einleitung}
\label{chap:einleitung}
\thispagestyle{fancy}

Die vorliegende B-Prüfung \texttt{B-DBI02-XX3-N01} hat die Entwicklung einer webbasierten Umfrageanwendung zum Thema \glqq{}\textbf{\textit{Umfrage zum Studium}}\grqq{} zum Ziel. Die Anwendung besteht aus mehreren \acs{PHP}-Skripten, die auf einem Apache-Webserver mit einer Maria\acs{DB}-Datenbank im Hintergrund laufen.

Für die Entwicklung wird die \acs{XAMPP}-Umgebung gewählt, da sie alle notwendigen Komponenten (\textit{Apache}, \textit{Maria\acs{DB}}/\textit{My\acs{SQL}}, \textit{\acs{PHP}}) in einem einfachen Paket bereitstellt und plattformunabhängig ist. Die vorliegende Dokumentation beschreibt die vollständige Einrichtung der Entwicklungsumgebung unter Windows sowie die Vorbereitung der Datenbank.

\section{Installation von \acs{XAMPP}}

\acs{XAMPP} ist ein kostenloses und quelloffenes Paket, das einen Apache-Webserver, eine Maria\acs{DB}-Datenbank sowie \acs{PHP} in einer einfachen Installation bündelt.

\subsection{Herunterladen von \acs{XAMPP}}

Die aktuelle Version von \acs{XAMPP} wird von der offiziellen Webseite heruntergeladen:

\begin{center}
    \textsl{\url{https://www.apachefriends.org/de/index.html}}
\end{center}

Für Windows wird die Installationsdatei im \texttt{.exe}-Format angeboten.

\subsection{Installation unter Windows}

Die Installation unter Windows erfolgt über die heruntergeladene \texttt{.exe}-Datei. Hierzu werden folgende Schritte durchgeführt:

\begin{enumerate}
    \item Die \texttt{.exe}-Datei wird mit Administratorrechten gestartet.
    \item Die Sicherheitsabfrage von Windows wird mit \glqq{}\textit{Ja}\grqq{} bestätigt.
    \item Bei der Auswahl der Komponenten werden folgende Module aktiviert:
          \begin{itemize}
              \item \textbf{Apache} – Webserver
              \item \textbf{My\acs{SQL}} (\textit{Maria\acs{DB}}) – Datenbankserver
              \item \textbf{\acs{PHP}} – Skriptsprache
              \item \textbf{phpMyAdmin} – grafische Datenbankverwaltung
          \end{itemize}
          Nicht benötigte Komponenten wie FileZilla, Mercury und Tomcat werden deaktiviert.
    \item Als Installationsordner wird \texttt{C:\textbackslash xampp} gewählt (\textit{oder ein alternativer Pfad ohne Leerzeichen}).
    \item Der Installationsassistent wird bis zum Abschluss durchlaufen.
    \item Nach der Installation wird die Option \glqq{}\textit{\acs{XAMPP} jetzt starten}\grqq{} aktiviert.
\end{enumerate}

\section{Starten der Dienste}

Nach erfolgreicher Installation müssen die benötigten Dienste gestartet werden.

\subsection{\acs{XAMPP} Control Panel öffnen}

Das \acs{XAMPP} Control Panel wird über das Startmenü unter \glqq{}\textit{XAMPP Control Panel}\grqq{} geöffnet.

\subsection{Starten von Apache und My\acs{SQL}}

Im Control Panel werden folgende Dienste gestartet:

\begin{itemize}
    \item \textbf{Apache} – Der Webserver, der \acs{PHP}-Seiten ausliefert.
    \item \textbf{My\acs{SQL}} – Der Datenbankserver (\textit{Maria\acs{DB}}), der die Umfragedaten speichert.
\end{itemize}

Nach dem Start leuchten die entsprechenden Buttons im Control Panel grün.

\section{Überprüfung der Umgebung}

Um sicherzustellen, dass alle Komponenten korrekt funktionieren, werden folgende Tests durchgeführt:

\subsection{Test des Apache-Webservers}

Der Browser wird geöffnet und die Adresse \textsl{\url{http://localhost/}} aufgerufen. Die \acs{XAMPP}-Startseite wird erfolgreich angezeigt.

\subsection{Test von \acs{PHP}}

Im Verzeichnis \texttt{\acs{htdocs}} (\textit{\acs{z. B.} \texttt{C:\textbackslash xampp\textbackslash htdocs\textbackslash}}) wird eine Datei \texttt{test.php} mit folgendem Inhalt erstellt:

\lstinputlisting[
    language=php, 
    caption={PHP test page}, 
    label={lst:php_test_page}
]{asset/code/Einleitung.tex}

Der Aufruf \textsl{\url{http://localhost/test.php}} zeigt die \acs{PHP}-Infoseite mit allen Konfigurationsdetails.

\subsection{Test von Maria\acs{DB} / phpMyAdmin}

Der Browser wird mit \textsl{\url{http://localhost/phpmyadmin/}} geöffnet. Die phpMyAdmin-Oberfläche wird erfolgreich geladen und zeigt die vorhandenen Datenbanken an.

\section{Projektverzeichnis vorbereiten}

Für die Entwicklung der Umfrageanwendung wird das folgende Verzeichnis im \texttt{\acs{htdocs}}-Ordner angelegt:

\begin{forest}
    for tree={
        font=\ttfamily,
        grow'=0,
        child anchor=west,
        parent anchor=south,
        anchor=west,
        calign=first,
        edge path={
            \noexpand\path [draw, \forestoption{edge}] (!u.south west) +(10pt,0) |- node[fill,inner sep=1.5pt] {} (.child anchor)\forestoption{edge label};
        },
        before typesetting nodes={
            if n=1{
                insert before={[,phantom,calign with current]}
            }{}
        },
        fit=band,
        before computing xy={l=22pt},
    }
    [\textbf{C:\textbackslash xampp\textbackslash htdocs\textbackslash survey}
        [\textbf{inc\textbackslash}]
    ]
\end{forest}

Dieses Verzeichnis dient als Wurzel der Anwendung. Alle \acs{PHP}-Dateien werden hier abgelegt, die datenbankrelevanten Dateien im Unterordner \texttt{inc/}.

Die Anwendung ist später unter der \acs{URL} \textsl{\url{http://localhost/survey/}} erreichbar.

\section{Anlegen der Datenbank}

Die Datenbank wird über phpMyAdmin erstellt:

\begin{enumerate}
    \item Aufruf von \textsl{\url{http://localhost/phpmyadmin/}}
    \item Klick auf \glqq{}\textit{Neu}\grqq{} in der linken Seitenleiste
    \item Datenbankname: \texttt{umfrage\_db}
    \item Zeichensatz: \texttt{utf8mb4\_general\_ci}
    \item Klick auf \glqq{}\textit{Anlegen}\grqq{}
\end{enumerate}

Die Datenbank ist nun bereit für die Tabellen, die im Rahmen der ersten Programmieraufgabe erstellt werden.

\section{Zusammenfassung der Umgebung}

Die folgende Tabelle fasst die wichtigsten Komponenten und deren Versionen zusammen:

\begin{center}
    \begin{tabular}{|l|l|}
        \hline
        \textbf{Komponente} & \textbf{Version} \\
        \hline
        Betriebssystem & Windows 11 \\
        \acs{XAMPP} & 8.2.4 \\
        Apache & 2.4.57 \\
        Maria\acs{DB} & 10.4.28 \\
        \acs{PHP} & 8.2.4 \\
        phpMyAdmin & 5.2.1 \\
        \hline
    \end{tabular}
\end{center}

% ==============================================================================
% =================================== Einleitung ================================
% ==============================================================================